% Gerado por regressao/06_tabelas_latex.R (projeto Modelagem). Não editar à mão.
\begin{landscape}
\begin{table}[htbp]
\caption{Robustez do ciclo eleitoral: modelo principal e testes selecionados}\label{tab:robustez_resumo}
\centering
\footnotesize
\setlength{\tabcolsep}{3pt}
\begin{tabular}{@{}>{\raggedright\arraybackslash}p{\dimexpr\linewidth-594pt\relax}>{\centering\arraybackslash}p{60pt}>{\centering\arraybackslash}p{60pt}>{\centering\arraybackslash}p{60pt}>{\centering\arraybackslash}p{60pt}>{\centering\arraybackslash}p{60pt}>{\centering\arraybackslash}p{60pt}>{\centering\arraybackslash}p{60pt}>{\centering\arraybackslash}p{60pt}>{\centering\arraybackslash}p{60pt}@{}}
\hline
\textbf{Variável} & \textbf{\hspace{0pt}Despesa total} & \textbf{\hspace{0pt}Saúde e Saneamento} & \textbf{\hspace{0pt}Educação e Cultura} & \textbf{\hspace{0pt}Habitação e Urbanismo} & \textbf{\hspace{0pt}Assistência Social} & \textbf{\hspace{0pt}Transporte} & \textbf{\hspace{0pt}Administração} & \textbf{\hspace{0pt}Agricultura} & \textbf{\hspace{0pt}Comunicações} \\ \hline
\multicolumn{10}{@{}l}{\textit{Modelo principal (erro de Driscoll-Kraay)}} \\
Ano eleitoral & 277,865\textsuperscript{***} & 53,163\textsuperscript{***} & 49,079\textsuperscript{*} & 126,814\textsuperscript{***} & 14,078\textsuperscript{***} & 41,741\textsuperscript{***} & 6,817 & 0,486 & $-$0,801\textsuperscript{***} \\
 & (102,741) & (10,703) & (29,440) & (31,902) & (4,274) & (10,585) & (10,968) & (2,026) & (0,218) \\
Ano pré-eleitoral & 162,806\textsuperscript{*} & 1,221 & 78,512\textsuperscript{**} & 34,083\textsuperscript{**} & 9,468\textsuperscript{**} & 8,342 & 10,347 & 5,168 & 0,656\textsuperscript{***} \\
 & (94,991) & (12,185) & (38,604) & (16,477) & (4,677) & (7,694) & (8,395) & (3,632) & (0,148) \\[6pt]
\multicolumn{10}{@{}l}{\textit{R5: especificação de Sakurai (2009) (erro agrupado por município)}} \\
Ano eleitoral & 202,329\textsuperscript{***} & 80,404\textsuperscript{***} & $-$40,640\textsuperscript{***} & 123,409\textsuperscript{***} & 10,825\textsuperscript{***} & 38,756\textsuperscript{***} & 9,046\textsuperscript{***} & $-$0,629 & $-$0,622 \\
 & (10,261) & (2,432) & (3,204) & (3,696) & (0,666) & (2,721) & (2,988) & (1,011) & (0,513) \\[6pt]
\multicolumn{10}{@{}l}{\textit{R12: erro-padrão agrupado por município}} \\
Ano eleitoral & 277,865\textsuperscript{***} & 53,163\textsuperscript{***} & 49,079\textsuperscript{***} & 126,814\textsuperscript{***} & 14,078\textsuperscript{***} & 41,741\textsuperscript{***} & 6,817\textsuperscript{*} & 0,486 & $-$0,801 \\
 & (12,099) & (2,749) & (3,742) & (4,270) & (0,788) & (3,159) & (3,560) & (1,165) & (0,579) \\
Ano pré-eleitoral & 162,806\textsuperscript{***} & 1,221 & 78,512\textsuperscript{***} & 34,083\textsuperscript{***} & 9,468\textsuperscript{***} & 8,342\textsuperscript{***} & 10,347\textsuperscript{***} & 5,168\textsuperscript{***} & 0,656\textsuperscript{***} \\
 & (5,513) & (1,549) & (1,981) & (2,115) & (0,492) & (2,020) & (1,972) & (0,758) & (0,224) \\[6pt]
\multicolumn{10}{@{}l}{\textit{R14: sem a dummy de pandemia (erro de Driscoll-Kraay)}} \\
Ano eleitoral & 259,057\textsuperscript{**} & 79,403\textsuperscript{***} & $-$10,566 & 134,645\textsuperscript{***} & 14,062\textsuperscript{***} & 41,614\textsuperscript{***} & 12,292 & 1,101 & $-$0,407 \\
 & (100,901) & (23,450) & (56,576) & (28,577) & (4,070) & (9,582) & (10,947) & (2,631) & (0,351) \\
Ano pré-eleitoral & 165,778\textsuperscript{*} & $-$2,926 & 87,930\textsuperscript{**} & 32,847\textsuperscript{*} & 9,470\textsuperscript{**} & 8,362 & 9,485 & 5,070 & 0,602\textsuperscript{***} \\
 & (95,259) & (14,674) & (41,429) & (16,773) & (4,673) & (7,739) & (8,329) & (3,638) & (0,184) \\
\hline
\end{tabular}
\fonte{Elaboração própria.}
\nota{Coeficientes do modelo A (efeito fixo de município), em R\$ de 2025 per capita. Modelo principal e R14: erros-padrão de Driscoll-Kraay; R5 e R12: erros-padrão agrupados por município. O R5 segue \textcite{sakurai2009}: sem ano pré-eleitoral e sem dummy de pandemia; o R14 retira só a dummy de pandemia, e 2020 volta a contar como ano eleitoral. Erros-padrão entre parênteses. \textsuperscript{***}p$<$0,01; \textsuperscript{**}p$<$0,05; \textsuperscript{*}p$<$0,1.}
\end{table}
\end{landscape}
